\section{引言}\label{sec:intro}

无人机为巡检、交通监测和灾害评估提供航拍图像。
如何从这些图像中提取有用信息，推动了人工辅助分析与自动目标检测的发展
\cite{ofli2016combining,pi2020convolutional}。
视觉问答提供了另一种交互方式：用户可以通过自然语言询问图像内容，例如场景中的车辆数量。
遥感视觉问答研究已建立这一任务，并提供了相应的评估基准\cite{lobry2020rsvqa}。
借助预训练视觉--语言模型（VLM）\cite{wang2024qwen2vl}，
边缘服务器能够结合航拍图像与问题进行处理。
然而，将这一方式用于无人机上行链路时，需要同时考虑通信与计算开销。
图像传输占用有限的无线资源，生成每个答案也需要模型推理。
因此，降低每题能耗需要同时考虑这两部分开销。

回答一个问题所需的信息并不一定包含完整图像。
例如，检测到的目标类别及其数量可能足以比较汽车与公交车的数量。
这类结构化信息的数据量较小，也可以通过简单规则处理。
但仅凭检测器输出，通常无法恢复漏检的目标。
图像保留的视觉细节可能有助于回答涉及漏检目标的问题，但也需要更多传输与接收端计算。
此外，更多视觉细节并不保证某个特定 VLM 对所有问题都能回答得更准确。
两种方式的选择取决于问题、已有检测信息，以及两条处理路径在通信链路上的实际表现。
有效的选择方法需要判断：何时较小的表示已经足够，何时传输图像值得付出额外开销。

语义通信与面向任务的通信为根据接收端任务的信息需求进行选择提供了基础
\cite{gunduz2023beyond,strinati2021sixg,kountouris2021semantics}。
针对无线视觉问答，MU-DeepSC 联合设计收发机以传输图像与问题特征
\cite{xie2022mudeepscwcl}。
面向目标的视觉问答通信则根据问题提取并排序目标框或场景图信息
\cite{liu2024gosg}。
面向查询的图像编码也利用查询选择用于图像重建的特征\cite{huang2026query}。
Park 和 Yoon 比较了不同任务下的视觉表示，并根据任务需求选择简单或丰富的语义信息
\cite{park2025transmit}。
这种任务层面的选择启发了本文进一步考察同一视觉问答工作负载中，每个问题应采用哪类证据与解码器。
Du 等利用 LLaVA 处理车载视觉问答，并根据用户关注内容确定图像区域及图像块传输的优先级
\cite{jiang2025lmmvn}。
该方法调整多模态模型的视觉输入；本文则研究结构化证据何时足以支持符号解码，从而避免 VLM 推理。
这些工作已经将任务感知与问题感知传输确立为已有设计思路。
在视觉推理方面，Prompt-RSVQA 将提取的视觉上下文转为文本，供语言模型使用
\cite{chappuis2022promptrsvqa}。
因此，以结构化视觉信息作为推理模型的另一类输入，也已有研究基础。

另一类研究关注语义处理的能耗。
绿色语义通信在空中辅助边缘网络及概率语义通信系统中同时考虑通信与计算开销
\cite{zheng2024aerialee,dai2025pscom}。
输入感知的 VLM 路由则根据输入及准确率与开销之间的权衡，在边缘与云端模型之间进行选择
\cite{sabanovic2026inarvl}。
在这些研究基础上，本文研究结构化航拍视觉问答中的一个具体决策：
传输检测信息并进行符号解码，还是传输图像并执行 VLM 推理。
该决策同时改变传输表示与回答问题所需的计算。
因此，需要结合问答准确率以及通信与计算的总能耗评估其价值，仅比较载荷大小并不充分。

为研究这一决策，本文提出一种包含发送端路由器和两条接收处理路径的证据级语义通信系统。
检测证据路径传输结构化检测结果，并使用轻量符号解码器；图像路径将图像传输至冻结的 VLM。
系统采用适用于航拍领域的检测器和常规信道编码，无需联合训练收发机与 VLM。
免训练规则根据问题类型选择路径。
学习型逐样本路由器进一步利用检测信息，预测哪条路径能够正确回答，并通过能耗价格控制准确率与能耗之间的权衡。
我们在涉及检测目标类别的结构化问题上评估这些选择。
研究目标是确定证据选择何时能够在不降低问答准确率的情况下提高能效，并刻画收益成立的条件。
本文的主要贡献如下：

\begin{itemize}
\item 将证据选择建模为传输载荷与接收端处理方式的联合选择。
能耗核算包括传输、机载检测与接收端推理，并区分无人机与边缘端的开销。
该建模明确了各条问答路径的代价，为能耗与准确率的比较提供依据。
\item 设计免训练类型规则和基于 MLP 的逐样本路由器。
条件性分析说明类型规则在何种条件下能够保持准确率并降低能耗。
引入能耗价格后，可调节学习型选择器的运行点。
\item 在两个航拍数据集上评估系统，主数据集覆盖 AWGN、Rayleigh 与 Rician 信道，
第二数据集采用 Rician 信道。
评估结合已记录结果的回放、GPU 功耗测量与传输开销建模。
在 VisDrone 的参考配置下，相较于速率自适应图像传输，类型规则将核算的平均每题能耗降低约 $55\%$，同时提高准确率。
在 DroneVehicle 上，验证集选择工作点的学习型路由以每题 $4.47$ J 达到 $74.60\%$ 准确率，
类型规则则为 $11.21$ J 和 $72.93\%$。
保留简单基线与匹配接收器比较，说明学习型选择何时有帮助，以及何时额外复杂度的收益有限。
这些结果针对系统合计能耗，不代表实测无人机电池收益。
\end{itemize}

本文其余部分安排如下：第~\paperref{sec:related}~节回顾相关工作；
第~\paperref{sec:system}~节介绍系统模型与能耗核算；
第~\paperref{sec:routing}~节介绍选择方法及其分析依据；
第~\paperref{sec:exp}~节报告实验并讨论其局限；第~\paperref{sec:conc}~节总结全文。
